\section{Peripheral eigenvalue group structure}

This section proves the peripheral spectrum structure of~\cite[Theorem~6.6]{Wolf2012Quantum} for irreducible Schwarz maps:
the peripheral eigenvalues form a finite cyclic group.
The trace-preserving refinement that the order divides the bond
dimension is proved in
Theorem~\ref{thm:peripheral_cyclic_group}.

\begin{lemma}[Peripheral eigenvectors are invertible]%
    \label{lem:isUnit_peripheral_eigenvector}
    \lean{MPSTensor.isUnit_peripheral_eigenvector}
    \leanok
    \uses{thm:ks_peripheral}
    Let $E$ be an irreducible unital Kraus map on $\MN{D}$ with a
    positive-definite adjoint fixed point.
    If $X \ne 0$ satisfies $E(X) = \mu X$ with $|\mu| = 1$,
    then $X$ is invertible.
\end{lemma}

\begin{proof}
    \leanok
    The KS equality gives $E(X^\dagger X) = X^\dagger X$, so $X^\dagger X$
    is a nonzero PSD fixed point. Irreducibility upgrades it to positive
    definite, hence invertible, so $\det(X) \ne 0$.
\end{proof}

\begin{lemma}[Peripheral eigenvalues: product closure]%
    \label{lem:peripheral_mul_closure}
    \lean{MPSTensor.peripheralEigenvalues_mul_mem_of_irreducible_unital_of_adjoint_fixedPoint}
    \leanok
    \uses{lem:isUnit_peripheral_eigenvector,thm:ks_peripheral}
    Let $E$ be an irreducible unital Kraus map on $\MN{D}$ with a
    positive-definite adjoint fixed point.
    If $\mu, \nu$ are peripheral eigenvalues of $E$, then so is $\mu\nu$.
\end{lemma}

\begin{proof}
    \leanok
    Take eigenvectors $X, Y$ for $\mu, \nu$. The KS equality at $X$
    puts $X$ in the multiplicative domain, so
    $E(YX) = E(Y)E(X) = (\nu\mu)(YX)$. Since $X, Y$ are units
    (Lemma~\ref{lem:isUnit_peripheral_eigenvector}), $YX \ne 0$ and
    $|\mu\nu| = 1$.
\end{proof}

\begin{theorem}[Peripheral eigenvalues: cyclic structure]%
    \label{thm:peripheral_cyclic_structure}
    \lean{PeripheralSpectrum.peripheral_eigenvalues_cyclic_structure}
    \leanok
    \uses{thm:ks_peripheral}
    Let $E$ be an irreducible unital Schwarz map on $\MN{D}$ with a
    positive-definite adjoint fixed point (trace-preservation is \emph{not}
    required).
    Then there exist $m \ge 1$ and a primitive $m$-th root of unity
    $\gamma$ such that the peripheral eigenvalues of $E$ are exactly
    $\{1, \gamma, \gamma^2, \ldots, \gamma^{m-1}\}$.
\end{theorem}

\begin{proof}
    \leanok
    The peripheral eigenvalues form a finite subset of $\C$,
    contain $1$, and are closed under multiplication and inversion,
    so they form a finite subgroup of $\C^\times$.
    Any finite subgroup of $\C^\times$ is cyclic.
    The generator $\gamma$ has order $m = |\text{peripheral spectrum}|$.
\end{proof}
